% !TeX root = ../papers/nc-p3.tex
\section{曲線の二点縮約の階数低下因子}
\label{sec:opposite-curve-divisor}
四辺の左右射影ファイバーが空か一点で、標準有限Frobenius候補の点グラフが
滑らかな整曲線成分 $C$ を持つとする。$C$ が全位置で射影三空間を張り、
標準四移動がこの成分を保つとする。点列に沿う埋込み $j_i:C\to\mathbb P(V_i^*)$
で $j_{i+4}=j_i\circ\sigma$ となるよう $\sigma$ を定める。
$L_i=j_i^*\mathcal O(1)$、$K=L_i L_{i+2}/L_{i+1}^2$ は前出の一定の多重差である。
\begin{theorem}\label{thm:opposite-curve-divisor}
$C\times C$ の点 $(p,q)$ で、位置0と2の二点 $j_0(p),j_2(q)$ により $w$ を縮約する。
階数が1以下になる軌跡は有効Cartier因子 $D$ で、その類は
\[
 \mathcal O(D)=\mathcal O(\Delta+\Gamma_{\sigma^{-1}})
       \otimes\operatorname{pr}_1^*K^{-1}\otimes\operatorname{pr}_2^*K^{-1}.
\]
ここで $\Delta$ は対角、$\Gamma$ はグラフである。$D$ は両側に有限平坦で次数2となる。
\end{theorem}
\begin{proof}
$H_1(p),H_3(p)$ を $j_1(p),j_3(\sigma^{-1}p)$ がそれぞれ消す超平面とし、
これらを第一曲線上の部分束と見る。$S=\operatorname{pr}_1^*L_0\otimes
\operatorname{pr}_2^*L_2$ とする。二点縮約は隣接点での消滅により
$H_3^*\to H_1\otimes S$ を与える（第一因子からの引戻しを略記する）。
その右核には $j_3(q)$、左核には $j_1(q)$ の制限が入る。
それぞれの三座標の共通零点は $q=\sigma^{-1}p$, $q=p$ のグラフで、スキームとして重複度1である。
実際、座標は一つの埋込みに対する点の座標差を局所的に生成し、
滑らかな曲線の対角イデアルは可逆である。従ってこの共通因子を除けば
右・左の核にそれぞれ直和局所因子となる直線束
\[
 A=\operatorname{pr}_2^*L_3^{-1}(\Gamma_{\sigma^{-1}}),\qquad
 B=\operatorname{pr}_2^*L_1^{-1}(\Delta)
\]
が入る。局所的な因子による除算後も縮約が零なのは曲面が整だからである。
そこで二点行列は $(H_3^*/A)\to B^\perp\otimes S$ という二階数束の間の写像となる。

固定した $p$ の上でこの写像が恒等的に階数1以下になることはない。
仮にそうなら、$j_2(C)$ を三行三列の行列空間に線型に写した像は
零行列を除いて階数1である。両端の射影像はいずれも射影平面を張る。
これは $j_2(C)$ が射影三空間を張り、$w(j_0(p))$ の端支持がともに3であるためである。
関数体上で行列を列と行の積に書くと、両側の三成分が張る関数空間は各三次元である。
その積は少なくとも五次元となる。実際、各空間で一つの滑らかな点の位数が相異なる基底を取り、
二つの三元整数集合の和集合（対ごとの和の集合）には五つ以上の位数が現れる。
ところが九成分は元の四座標の線型結合なので四次元以下であり矛盾する。
$q$ を固定しても同様である。

よって二行二列の行列式は零でなく、階数低下を切る。
$\det H_1=L_1^{-1},\ \det H_3=(\sigma^{-1})^*L_3^{-1}$ より行列式の束は
\[
 S^2\otimes \det H_1\otimes\det H_3\otimes A\otimes B.
\]
$L_{-1}=(\sigma^{-1})^*L_3$ と束の三階差を用いれば主張の類となる。
どのファイバーも因子に含まれず、制限次数は2である。従って固有準有限射として有限、
純曲線から滑らかな曲線へ台の落ちる成分がないので平坦である。
\end{proof}

\begin{corollary}\label{cor:curve-divisor-monodromy}
次のモノドロミー制限が成り立つ。
\begin{enumerate}
\item 種数 $g\ge2$ で $\sigma=1$、または $\sigma\ne1$ で固定点が $2g-2$ 未満なら
$D=\Delta+\Gamma_{\sigma^{-1}}$ かつ $K=1$ である。
\item $g\ge2$ で $\sigma$ が対合なら、$K\ne1$ が可能なのは商が楕円曲線のときだけである。
この場合も $K^2=1$ であり、商写像 $f:C\to E$ に対して $D$ は
$f(q)-f(p)=e$、$e\in E[2]$ の一ファイバーとなる。
\item $C$ が楕円曲線で $\sigma$ が平行移動なら $K^2=1$、
$D$ は重複を許す二つの平行移動グラフの和である。
一般の楕円四次成分でも常に $K^6=1$ である。
\end{enumerate}
\end{corollary}
\begin{proof}
対角とグラフの自己交叉は $2-2g$ である。相異なる場合の交叉は固定点の個数である。
非恒等の有限位数自己同型の固定点は標数零で単純である。
交叉数が負なら該当成分が $D$ に入るので、一方を引き、残りの自己交叉も負であることを使う。
$\sigma=1$ でも二度引けば同じである。類を一方のファイバーに制限すると $K=1$ を得る。

対合の商の種数を $b$ とすると固定点数は $2g+2-4b$ である。
$b\ge2$ なら第一項に帰着する。$b=0$ なら二重被覆の束を $H=f^*\mathcal O(1)$ として、
二グラフの和の束は $H\boxtimes H$ である。$D$ は縦横成分を持たないので
$h^0(H/K)\ge2$ が必要である。種数2以上では次数2の鉛筆は唯一である。
実際、二つあれば両射の像は $\mathbb P^1\times\mathbb P^1$ の $(2,2)$ 曲線に双有理で
種数が1以下になるか、同じ二重被覆を通る。従って $K=1$ である。

$b=1$ なら差の射 $d(p,q)=f(q)-f(p)$ の零ファイバーが二グラフの和である。
一般ファイバーは二つの二次関数体の合成に対応して整となる。
実際、分岐因子は非空で有限であり、その一般の平行移動とは一致しない。
Stein分解で全ファイバーが連結となり、全てのファイバーは一般点で分離的な
二重被覆の積として被約である（Cartier因子の純性を使う）。
$D$ はファイバーと数値同値なので各成分が $d$ で点に写る。
各ファイバーの成分の交叉行列の核はファイバー全体の重複度列が張る。
実際、最大係数の成分でファイバーの定数倍を引いて交叉数零を使えば、
隣接成分も同じ係数となり連結性で伝わる。
ゆえに $D$ は丸ごとの一ファイバーである。両座標ファイバーへの制限から
$K$ は $E$ 上の互いに逆な二つの次数零束のどちらの引戻しでもある。
$f^*$ は単射である。非自明な束が自明になるならノルムで二乗が自明となり、
対応する非自明なエタール二重被覆を $f$ が経由するが、分岐に反する。
従ってこの二束は同じ二捩れである。

楕円曲線の平行移動の場合は、二グラフはいずれも $q-p$ のファイバーであり、
数値同値を用いれば $D$ も二ファイバーの和となる。対角に沿って次数零の束の差を制限すると
$K^2=1$ である。他の楕円自己同型では定理\ref{thm:point-graph-curve-degree}の
$uK=K$ により位数は2または3を割る。従って楕円四次ではどの場合も6で消える。
\end{proof}

\begin{corollary}\label{cor:septic-canonical-monodromy}
定理\ref{thm:septic-bielliptic-forms}の全ての滑らかな七次窓で、標準四移動が
点曲線に誘導する自己同型は $h\mapsto-h$ である。
商を $f:C\to E$、順序付き二点の像を $P,Q$ とすれば
\[
 D=\{(a,b):f(b)-f(a)=P-Q\},\qquad P-Q\in E[2]\setminus\{0\}.
\]
\end{corollary}
\begin{proof}
偏極が二つの標準射影を交互に繰り返すので、二階差は $\mathcal O(2p-2q)$ またはその逆である。
これは非自明な二捩れだから、前系で標準四移動は恒等になれない。
八係数の一般点では、順序付き二点を保つ曲線自己同型は恒等と $h$ の符号変更だけである。
実際、定理\ref{thm:septic-bielliptic-forms}の二つの倍率不変量が有限交換を排除し、
残る倍率は共通倍率と $h$ の符号だけである。標準四移動は両偏極を保ち二点を個別に保つので、
一般点で表示の対合に一致する。この一致は全滑らかな係数に及ぶ。
実際、一定のFrobenius次元を持つこの開集合上でtraceを局所的に取り、
完全積対合の逆行列から四移動の射影作用が正則に定まる。traceの取り替えは射影作用を変えない。
滑らかな完全交叉族は平坦で全空間に一般ファイバーが稠密であり、二射の一致は閉条件である。
両点は $h=0$ の分岐点なので $K=f^*\mathcal O(P-Q)$、$f^*$ の単射性と前系から因子の式を得る。
\end{proof}
